% Part: proof-theory
% Chapter: proof-search
% Section: tableaux

\documentclass[../../../include/open-logic-section]{subfiles}

\begin{document}

\olfileid{pt}{ps}{tab}

\olsection{टॅब्लो}

टॅब्लो पद्धती या क्रमवर्ती-कलनाशी जवळून
संबंधित !!{proof} पद्धती आहेत. हाताने किंवा
संगणकीय अंमलबजावणीत सिद्धता-शोधासाठी त्या विशेष
सोयीच्या आहेत. टॅब्लोत क्रमवर्तींचे वृक्ष न बांधता
!!a{proof} ही !!{formula}s चा वृक्ष असते.
पण नैसर्गिक निगमनापेक्षा वेगळे म्हणजे त्यांचे नियम
क्रमवर्ती-कलनाच्या नियमांशी जुळतात. टॅब्लो
!!{proof} म्हणजे मूलतः प्रतिमानाचा स्पष्ट शोध.
म्हणून त्यांना कधी \emph{चिन्हार्थविषयक} टॅब्लो,
तर कधी \emph{सत्यता-वृक्ष} म्हणतात.

\emph{चिन्हांकित} टॅब्लो हा चिन्हांकित
!!{formula}s चा वृक्ष आहे. चिन्हांकित सूत्राचे रूप
$\sFmla{\True}{!A}$ किंवा~$\sFmla{\False}{!A}$
असे असू शकते. वृक्ष खालच्या दिशेने वाढतात.
सुरुवातीला आपल्याला काय !!{prove} करायचे ते
दर्शवणारी !!{formula}s असतात. उदाहरणार्थ,
$!A, !B \Sequent !C, !D$ !!{prove} करण्यासाठी
वृक्षाची सुरुवात
$\sFmla{\True}{!A}, \sFmla{\True}{!B},
\sFmla{\False}{!C}, \sFmla{\False}{!D}$ अशी
होईल. ~$!A$ आणि~$!B$ सत्य, पण~$!C$ आणि~$!D$
असत्य ठरवणारी !!^a{structure} म्हणजे
$!A, !B \Sequent !C, !D$ अवैध असल्याचे
दाखवणारी !!a{structure}.

टॅब्लोचे तळाचे पानशिखर शाखा-विस्तार
नियमांनी वाढवता येते. हे नियम त्याच शाखेत नवी
शिखरे जोडतात किंवा खाली दोन भावंड शिखरे जोडून
शाखा दोन नव्या शाखांत विभागतात. त्याच शाखेतील
त्या शिखराच्या वरच्या कोणत्याही चिन्हांकित
!!{formula} ला नियम लागू करता येतात. अभिजात
टॅब्लो पद्धत~\Log{Tc} साठीचे शाखा-विस्तार नियम
\olref{tab:Tc} मध्ये दिले आहेत.

\subfile{rules-Tc}

नियम म्हणजे~\Log{G3c} चेच नियम उलटे
दाखवले आहेत. चिन्हे~$\True$ आणि~$\False$
अनुक्रमे क्रमवर्तीच्या डाव्या किंवा उजव्या बाजूला
!!a{formula} मुख्य आहे हे दर्शवतात.

टॅब्लोची शाखा~$\sFmla{\True}{!A}$ आणि
$\sFmla{\False}{!A}$ ही दोन्ही सूत्रे धरत असेल,
तर ती \emph{बंद} आहे. प्रत्येक शाखा बंद असेल,
तर संपूर्ण टॅब्लो बंद म्हणतो. उदाहरणार्थ,
$!A \lor (!B \land !C) \Sequent (!A \lor !B) \land
(!A \lor !C)$ साठी हा बंद टॅब्लो आहे; म्हणजे
$\sFmla{\True}{!A \lor (!B \land !C)}$ आणि
$\sFmla{\False}{(!A \lor !B) \land (!A \lor !C)}$
यांपासून सुरुवात होते:
\begin{oltableau}
  [\sFmla{\True}{\formula{A} \lor (\formula{B} \land \formula{C})},
    just=\TAss, checked
    [\sFmla{\False}{(\formula{A} \lor \formula{B}) \land
        (\formula{A} \lor \formula{C})}, just=\TAss, checked
      [\sFmla{\True}{\formula{A}}, just={\TRule{\True}{\lor}[1]}
        [\sFmla{\False}{\formula{A} \lor \formula{B}},
          just={\TRule{\False}{\land}[2]}, checked
          [\sFmla{\False}{\formula{A}},
            just={\TRule{\False}{\lor}[4]}
            [\sFmla{\False}{\formula{B}},
              just={\TRule{\False}{\lor}[4]}, close]
          ]
        ]
        [\sFmla{\False}{\formula{A} \lor \formula{C}},
          just={\TRule{\False}{\land}[2]}, checked
          [\sFmla{\False}{\formula{A}},
            just={\TRule{\False}{\lor}[4]}
            [\sFmla{\False}{\formula{C}},
              just={\TRule{\False}{\lor}[4]}, close]
          ]
        ]
      ]
      [\sFmla{\True}{\formula{B} \land \formula{C}},
        just={\TRule{\True}{\lor}[1]}, checked
        [\sFmla{\False}{\formula{A} \lor \formula{B}},
          just={\TRule{\False}{\land}[2]}, checked
          [\sFmla{\True}{\formula{B}},
            just={\TRule{\True}{\land}[3]}, move by=2
            [\sFmla{\True}{\formula{C}},
              just={\TRule{\True}{\land}[3]}
              [\sFmla{\False}{\formula{A}},
                just={\TRule{\False}{\lor}[4]}
                [\sFmla{\False}{\formula{B}},
                  just={\TRule{\False}{\lor}[4]},close
                ]
              ]
            ]
          ]
        ]
        [\sFmla{\False}{\formula{A} \lor \formula{C}},
          just={\TRule{\False}{\land}[2]}, checked
          [\sFmla{\True}{\formula{B}},
            just={\TRule{\True}{\land}[3]}, move by=2
              [\sFmla{\True}{\formula{C}},
                just={\TRule{\True}{\land}[3]}
                [\sFmla{\False}{\formula{A}},
                  just={\TRule{\False}{\lor}[4]}
                  [\sFmla{\False}{\formula{C}},
                  just={\TRule{\False}{\lor}[4]},close
                ]
              ]
            ]
          ]
        ]
      ]
    ]
  ]
\end{oltableau}

खूणा दर्शवतात की संबंधित नियम त्या
!!a{formula} च्या खालच्या सर्व शाखांवर लागू
केला आहे. ~$\otimes$ हे चिन्ह शाखा बंद असल्याचे
दर्शवते. नियमानंतरच्या ओळ-क्रमांकांवरून नियम
कोणत्या चिन्हांकित !!{formula}s ला लागू केला
हे कळते (म्हणजे दर्शवलेल्या ओळीवरील त्याच
शाखेतील !!{formula} ला).

टॅब्लोमध्ये सिद्धता-शोध सोपा असतो.
टॅब्लो टप्प्याटप्प्याने तयार करतो. ~$0$ व्या
टप्प्यावर सुरुवातीची !!{formula}s वर लिहा.
~$n+1$ व्या टप्प्यावर प्रत्येक पानशिखरासाठी
पुढील कृती करा: अद्याप खूण न केलेले सर्वांत
पहिले (सर्वांत वरचे) चिन्हांकित सूत्र शोधा.
त्याला शाखा-विस्तार नियम लावा. तो चिन्हांकित
सूत्र असलेल्या प्रत्येक शाखेवर नियम लागू झाल्यावर
त्याला खूण करा. (वरचे उदाहरण याच अल्गोरिदमने
बनवले आहे.)

~$\sFmla{\True}{\lforall}$ आणि
$\sFmla{\False}{\lexists}$ या रूपातील
चिन्हांकित !!{formula}s वर खूण कधीच करत नाही;
म्हणून त्यांना विशेष पद्धतीने हाताळावे लागते.
ती असतील, तर~$\TRule{\True}{\forall}$ आणि
$\TRule{\False}{\lexists}$ हे नियम लागू पडणाऱ्या
सर्व पदांसाठी अखेरीस लावले जातील याची खात्री
करावी लागते. उदाहरणार्थ, प्रत्येक टप्प्यावर,
त्या रूपातील नसलेले सर्वांत वरचे चिन्हांकित
!!{formula} हाताळल्यानंतर, प्रत्येक शाखेवरील
अशा सर्व सूत्रांना हे नियम लावता येतील.
वापरायचे पद~$t$ हे शाखेत आधीपासून असलेल्या
!!{constant}s पासून (एकही !!{constant} नसेल,
तर~$\Obj c_0$ पासून) आणि सुरुवातीच्या
!!{formula}s मधील !!{function}s पासून
बनवता येणारे पहिले पद~$t$ असते; तसेच अनुरूप
चिन्हांकित !!{formula}
$\sFmla{\True}{!B(t)}$ किंवा~$\sFmla{\False}{!B(t)}$
त्या शाखेत आधीपासून नसते.

या सिद्धता-शोधातून बंद टॅब्लो मिळाला नाही,
तर त्यात सांत खुली शाखा असते, जिच्यावरील सर्व
अनणुरूप !!{formula}s वर खूण झालेली असते;
किंवा शोध अनंत, सांत-शाखित वृक्ष तयार करतो,
ज्यात अनंत शाखा असतेच. \olref[cpl]{sec}
प्रमाणे, अशी !!a{structure}~$\Struct{M}$ बांधू की
शाखेवर~$\sFmla{\True}{!A}$ असेल, तर
$\Sat{M}{!A}$, आणि~$\sFmla{\False}{!A}$
असेल, तर~$\Sat/{M}{!A}$. (येथे
$\Theta = \Setabs{!A}{\sFmla{\True}{!A} \text{ शाखेवर}}$
आणि~$\Xi = \Setabs{!A}{\sFmla{\False}{!A} \text{ शाखेवर}}$
घ्या.)

या पद्धतीने तयार झालेले टॅब्लो प्रत्येक
शिखराला क्रमवर्ती देऊन~\Log{G3c} मधील सिद्धतांमध्ये
सहज रूपांतरित करता येतात. सुरुवातीची
सूत्रे ही~$\Gamma \Sequent \Delta$
क्रमवर्तीशी संबंधित असतील, तर त्या प्रत्येक
सूत्राला~$\Gamma \Sequent \Delta$ ही खूण द्या.
~$\sFmla{\True}{!A}$ या चिन्हांकित !!{formula}
वर शाखा-विस्तार नियम लावून शाखा वाढवली, तर
पानाला~$!A, \Pi \Sequent \Lambda$ ही खूण द्या;
~$\sFmla{\False}{!A}$ या चिन्हांकित !!{formula}
वर नियम लावला,
तर पानाला~$\Pi \Sequent \Lambda, !A$ ही खूण द्या.
टॅब्लोत~\TRule{\True}{\star} नियम वापरला,
तर क्रमवर्तीवर~\LeftR{\star} नियम लावा
(मुख्य सूत्र~$!A$);~\TRule{\False}{\star}
वापरला असेल, तर~\RightR{\star} लावा.
नियमाची आधारविधाने ही टॅब्लोत नव्याने
जोडलेल्या अनुरूप शिखरांना दिलेल्या क्रमवर्ती आहेत.
~$\sFmla{\True}{!A}$ आणि~$\sFmla{\False}{!A}$
यांनी टॅब्लो बंद झाल्यावर अंतिम पानशिखराची
खूण~$!A, \Pi \Sequent \Lambda, !A$
या रूपातील क्रमवर्ती असेल. टॅब्लोतील काही सलग
शिखरांना तीच क्रमवर्ती दिली जाईल; तिच्या
पुनरावृत्ती काढा. उदाहरणार्थ, वरच्या टॅब्लोचे
पुढील !!{proof} मध्ये रूपांतर होते:
\[\small
\Axiom$!A \fCenter !A, !B$
\RightLabel{\RightR{\lor}}
\UnaryInf$!A \fCenter !A \lor !B$
\Axiom$!A \fCenter !A, !C$
\RightLabel{\RightR{\lor}}
\UnaryInf$!A \fCenter !A \lor !C$
\RightLabel{\RightR{\land}}
\BinaryInf$!A \fCenter (!A \lor !B) \land (!A \lor !C)$
\Axiom$!B, !C \fCenter !A, !B$
\RightLabel{\RightR{\lor}}
\UnaryInf$!B, !C \fCenter !A \lor !B$
\RightLabel{\LeftR{\land}}
\UnaryInf$!B \land !C \fCenter !A \lor !B$
\Axiom$!B, !C \fCenter !A, !C$
\RightLabel{\RightR{\lor}}
\UnaryInf$!B, !C \fCenter !A \lor !C$
\RightLabel{\LeftR{\land}}
\UnaryInf$!B \land !C \fCenter !A \lor !C$
\RightLabel{\RightR{\land}}
\BinaryInf$!B \land !C \fCenter (!A \lor !B) \land
(!A \lor !C)$
\RightLabel{\LeftR{\lor}}
\BinaryInf$!A \lor (!B \land !C) \fCenter (!A \lor !B) \land
(!A \lor !C)$
\DisplayProof
\]

\end{document}
